 



\section{GTK3 UI Layout and User Interaction}

\subsection{Application Context and Initialization}
The initialization phase of the GTK3 client application is a critical foundational step that establishes the software environment necessary to interface with the remote embedded system. This project presents the design and implementation of a robust Measurement Network Gateway, a high-performance userspace C application deployed on a Xilinx Zynq-7000 SoC (ZedBoard). The remote client features a graphical platform utilized to continuously monitor real-time data streams and dynamically configure the distant sensor devices over the TCP/IP network. To manage this, the graphical application relies on a centralized context structure (often defined as ctx), which holds pointers to all the critical GTK widgets—such as labels, entry fields, buttons, and drawing areas—as well as the underlying data buffers required for Cairo graph rendering.

Initialization begins with the \texttt{create\_gui()} function, which serves as the primary constructor for the user interface. During this phase, the application allocates memory for the complex data structures required to hold incoming telemetry, setting the initial states for minimum and maximum graph values, and clearing the arrays to prevent graphical artifacting. Furthermore, because this client application must safely interact with network sockets, the initialization phase is responsible for setting up the necessary threading primitives, including mutexes, to ensure that the graphical thread does not conflict with the background data-receiving thread. The system architecture employs a concurrent Producer-Consumer model utilizing POSIX threads (pthreads).

The main desktop window is instantiated using gtk\_window\_new(GTK\_WINDOW\_TOPLEVEL). The title is explicitly set to "Measurement Network Gateway - Real-Time Data Monitor" to provide clear context to the user. A default resolution of 1500x950 pixels is enforced via gtk\_window\_set\_default\_size(), ensuring that the dashboard launches with ample screen real estate to render highly detailed, dual-axis Cairo graphs without immediate user resizing. Finally, the window is bound to a fundamental destruction callback (on\_window\_destroy) via g\_signal\_connect, ensuring that when the user closes the application, all associated TCP sockets are gracefully terminated, memory is freed, and the GTK main loop is cleanly exited.


\subsection{Top-level Window and Main Layout}
GTK3 fundamentally differs from absolute-positioning UI frameworks by utilizing a dynamic, responsive "box-packing" model. The top-level window acts as the blank canvas, but it can only contain a single root widget. Therefore, the core of the dashboard's architecture relies on a master vertical box, vbox\_main, instantiated with gtk\_box\_new(GTK\_ORIENTATION\_VERTICAL, 10). This master container acts as the structural spine of the application, stacking all subsequent sub-panels, frames, and operational zones from the top of the window to the bottom.

To create a visually cohesive and intuitive user experience, the interface heavily leverages GtkFrame widgets. These frames draw subtle borders around logically grouped components and provide embedded text headers, such as "Connection", "Sensor Frequency Set(Hz)", and "Sensor Data (Live)". Inside these frames, horizontal boxes (hbox\_conn, hbox\_config) are packed. This allows widgets like the server IP entry (gtk\_entry\_new) and the connection buttons to sit neatly side-by-side in a single row. This layout strategy is entirely responsive; if the user maximizes the application window, the box-packing algorithm automatically redistributes the free space, keeping the control panels anchored and proportioned correctly without distorting the text or buttons.

A pivotal structural component of the main layout is the GtkPaned widget, which splits the lower section of the window horizontally. The gtk\_paned\_new(GTK\_ORIENTATION\_HORIZONTAL) function creates a draggable divider (set to an initial position of 350 pixels via gtk\_paned\_set\_position). This empowers the user to dynamically adjust the visual ratio between the raw numerical statistics on the left and the graphical plotting area on the right. This separation of concerns within the UI mirrors the backend architecture; it cleanly isolates the control logic (top), the raw data readout (left), and the analytical visualization (right), resulting in an industrial-grade dashboard that is highly legible and operationally efficient.


\subsection{Left Side: Drawing Area + Statistics}
The left side of the dashboard's paned layout is dedicated to delivering immediate, raw, and high-precision statistical readouts from the embedded gateway. The primary objective is to acquire high-precision data from a heterogeneous suite of sensors including a MAX31865 RTD-to-Digital converter for temperature, analog inputs (ADC), and digital I/O (switches and push buttons). To present this continuous stream of data cleanly, the left pane utilizes a vertical box (vbox\_sensors) acting as a tabular column.

Each individual sensor is granted its own horizontal row (hbox\_temp, hbox\_adc0, etc.) containing three distinct GtkLabelwidgets. The first label utilizes GTK's markup language (gtk\_label\_set\_markup) to render the sensor's name in bold text, immediately drawing the operator's eye. A fixed width is applied to this title label using gtk\_widget\_set\_size\_request to ensure that all subsequent data labels align perfectly into a uniform, readable column, much like a spreadsheet. The second label in the row is dedicated to displaying the processed numerical value (e.g., the temperature in Celsius or the ADC voltage), while the third label, aligned strictly to the right edge (GTK\_ALIGN\_END), displays the corresponding microsecond timestamp of that specific reading.

While the exact layout code dictates placing the raw statistics on the left and the Cairo GtkDrawingArea on the right side of the paned divider, they function as a unified data presentation layer. The drawing area acts as the visual counterpart to the statistics. It maps the raw timestamps and normalized integer values into a continuous sliding-window graph. By keeping the live numeric labels (statistics) visibly anchored on the left while the complex, CPU-intensive Cairo graphics render the historical trends on the right, the user receives both immediate situational awareness (the current exact temperature) and analytical context (the thermal trend over the last 10 seconds) simultaneously. This fulfills the primary objective of providing a fault-tolerant foundation for real-time remote monitoring


\subsection{Right Side: Scrollable Control Panel}
The control panel components of the UI—distributed across the top frames and the right side of the graph—are responsible for translating user intent into TCP payloads destined for the remote ZedBoard. The dashboard must safely allow remote users to dynamically adjust sampling rates or pause the stream. To achieve this safely, the UI employs highly constrained input widgets. For frequency configuration, GtkSpinButton widgets are utilized (gtk\_spin\_button\_new\_with\_range). By explicitly setting the allowable range from 10Hz to 10,000Hz, the UI mathematically prevents the operator from submitting an invalid polling rate that could crash the embedded kernel driver or overwhelm the network socket.

Each sensor (Temperature, ADC0, ADC1, Switch) features its own independent spin button paired with a dedicated "Set" button. The configure command parses space-delimited string payloads to dynamically update the independent polling frequencies of the ADC, temperature, and switch components on the fly. When a user clicks a "Set" button, the GTK application extracts the double-precision value from the spin button, formats it into the custom "configure" command string, and dispatches it over the network.

Crucially, the control panel implements strict state-management logic to ensure operational safety. Upon application launch, buttons such as "Start", "Stop", "Disconnect", and the various "Set" configuration buttons are deliberately greyed out and disabled using gtk\_widget\_set\_sensitive(widget, FALSE). They only become active once the "Connect" button successfully establishes a valid TCP socket with the server IP address provided in the GtkEntry field. On the right side of the graph, a localized control panel utilizes GtkComboBoxText dropdown menus. These allow the user to instantly select which specific sensor streams are mapped to the red and blue plotting axes of the Cairo drawing area, granting powerful, real-time customization of the data visualization without needing to restart the stream.


\subsection{Bottom: Command Terminal UI}
To provide deep diagnostic visibility into the network communication between the PC client and the embedded Linux gateway, the bottom of the dashboard layout features a dedicated command terminal UI. In complex industrial deployments, simply seeing a graph is often insufficient for debugging connection drops, packet corruption, or hardware fault flags originating from the remote sensors. The terminal acts as a transparent window into the underlying TCP/IP protocol, providing the operator with a raw, chronological feed of system events and data payloads.

This terminal is constructed using a GtkTextView widget embedded within a GtkScrolledWindow. The scrolled window is explicitly configured to handle vertical overflow, meaning that as thousands of lines of log data are generated, the view automatically provides a scrollbar rather than allowing the UI layout to expand off the physical screen. The text view itself is bound to a GtkTextBuffer, which holds the actual string data. To ensure the terminal acts purely as a diagnostic log and not an accidental input field, the widget is strictly locked using gtk\_text\_view\_set\_editable(FALSE).

As the background networking thread receives TCP packets, it unpacks the custom command protocol. If the client software evaluates the least significant bit (Bit 0) and detects a hardware fault flag, or if the connection status changes, these events are appended to the text buffer. A utility function is utilized to automatically move the text cursor (GtkTextIter) to the very end of the buffer after every insertion, forcing the scrolled window to automatically scroll down. This ensures the user is always looking at the absolute latest network event, creating a highly effective, live-scrolling diagnostic terminal that complements the graphical plotting area above it.


\subsection{Callbacks and Periodic Refresh}
The GTK3 framework relies on an event-driven architecture, meaning the application idles in a main loop (gtk\_main()) until an internal timer fires or a user interacts with a widget. This interaction is handled entirely through callbacks and signal routing. Every interactive widget in the dashboard is mapped to a specific C function using g\_signal\_connect. For example, clicking the "Connect" button triggers the on\_connect\_clicked callback, which immediately attempts to open a non-blocking TCP socket to the IP address specified in the UI.

Because the GTK main loop must never be blocked (otherwise the UI will freeze and the application will become unresponsive to the OS), the continuous rendering of the Cairo graph is managed via a periodic refresh timer. A function such as g\_timeout\_add() is initialized to fire every ~33 milliseconds (yielding approximately 30 frames per second). When this timer ticks, it triggers a callback that updates the numerical GtkLabel statistics and issues a gtk\_widget\_queue\_draw() command to the GtkDrawingArea. This queues an expose event, forcing GTK to invoke the on\_draw\_graph function, which subsequently recalculates the sliding time window and redraws the grid, axes, and sensor lines.

Thread synchronization is paramount within these callbacks. The background networking thread is constantly listening to the TCP socket and pushing new data points into the client-side buffers. Since the GTK UI thread and the networking thread run concurrently, accessing the graph data arrays simultaneously would cause fatal race conditions or corrupted screen renders. Therefore, whenever the on\_draw\_graph callback is fired by the refresh timer, it first executes a pthread\_mutex\_lock() on the shared context structure. It reads the historical data, plots the pixels, and then immediately calls pthread\_mutex\_unlock(). This guarantees atomic access to the data, ensuring a perfectly smooth, tear-free graphical user experience.
